\section{Network Namespace Separato per la
Verifica}\label{network-namespace-separato-per-la-verifica}

In BlueAgent, abbiamo introdotto OVN per semplificare la gestione di due
reti con stesso gateway.

Per evitare di aggiungere la complessità computazionale di OVN, si può
pensare di creare un network namespace separato al momento della
verifica.

All'interno di questo nuovo namespace, chiamato \(\text{verify-}k\) si
crea la rete di verifica composta dal solo bridge L2 e si assegna a
questo bridge il gateway originale.

I container duplicati dovranno poi essere collegati al nuovo bridge.

Per consentire l'accesso alla rete WAN bisogna usare una coppia di
\passthrough{\lstinline!veth!}, per collegare il namespace
\(\text{verify-}k\) al namespace root.

In questo modo i due bridge con stesso ip gateway possono coesistere
poiché risiedono in namespace differenti, ci pensa il kernel ad
effettuare il routing dei pacchetti verso le giuste interfacce nei
differenti namespace.

Tuttavia l'implementazione non è immediata e l'uso di Incus non si sposa
bene con la creazione di network namespace separati dato che non sa come
tracciarne l'esistenza. Dunque bisogna gestire manualmente le interfacce
di rete, aggiungendo una notevole complessità. Questo perché il motore
di Incus esegue sul namespace root, e non è progettato per gestire
interfacce dentro altri namespace. Bisogna dunque raggirare Incus, ed
istanziare e gestire in BlueAgent coppie di \emph{veth} che gestiscano
questo disallineamento.

Dal lato del codice aumenta dunque la complessità: invece di fare
richieste verso l'API di Incus per gestire le interfacce, bisogna usare
vari comandi e tracciare il ciclo di vita delle interfacce usate.
